%%
%% A true multi-typed WGSL kingdom: every gene a legal GPU program by its types,
%% applied to UltraDim's k-means centroid update.
%%
%% Domain paper (4) of the set; the shared material (the schema, the master symbol
%% table, the geneframe, GPU-resident evolution, HFF) is in the core paper
%% \cite{morgan_core}. Every number here has a source file named beside it.
%% Figures are drawn by pgfplots from wgsl_multityped_gep_data/*.tsv.
%%
\documentclass[sigconf,nonacm]{acmart}

\usepackage{amsmath}
\usepackage{booktabs}
\usepackage{tabularx}
\usepackage{pgfplots}
\usepackage{pgfplotstable}
\pgfplotsset{compat=1.18}

\usepackage{xurl}
\urlstyle{tt}
\newcommand{\code}[1]{\nolinkurl{#1}}
\newcommand{\src}[1]{{\footnotesize\nolinkurl{#1}}}

\begin{document}

\title{A True Multi-typed WGSL Kingdom}
\subtitle{Every gene a legal GPU program by its types, applied to UltraDim's
k-means centroid update}

\author{Andrew James Morgan}
\email{andrew@gamakon.ai}
\affiliation{%
  \institution{Gamakon Ltd.}
  \city{London}
  \country{United Kingdom}
}
\renewcommand{\shortauthors}{Morgan}

\begin{abstract}
We describe a kingdom of multi-typed Gene Expression Programming (GEP) whose
phenotypes are WGSL compute kernels, and in which every gene, at every position
and after every operator, expresses a legal kernel. Legality is not checked,
filtered or repaired: each rule of a legal GPU program is carried by a type of
the master symbol table, and a symbol exists only in the types where its use is
legal. The rules carried include an index in range of the buffer it indexes, a
thread that stores only to a cell it owns, a barrier only in uniform control, a
loop whose trip count is bounded by the problem's sizes, a workgroup pool no
larger than the tree names, a data-indexed store whose slot is assigned by the
compiler and never by the gene, and a scratch region that is read-only or
write-only within one dispatch. Python generators emit the tables and refuse to
write one unless witnesses, run over the emitted text, prove these properties.
We apply the kingdom to UltraDim's k-means centroid update, with UltraDim's own
bindings, so every individual is a drop-in. On an Apple M3 Max under Metal, the
best individual of one evolution run gives A7's output bit for bit on all 8
correctness cases and runs at 0.0274 of A7's time (11 of 11 interleaved pairs
faster). Almost all of that is the hand-written counting-sort form that the
population started from (0.0278 of A7). Evolution added 1.6\,\% over that form
(0.9839, 11 of 11 pairs). Most of that gain comes from removing a Kahan
compensation that, on Metal, changes no bit. We report the measurement faults
we found and fixed, the controls, and the open work.
\end{abstract}

\keywords{gene expression programming, strongly typed genetic programming, GPU,
WGSL, program legality by construction, genetic improvement, k-means}

\maketitle

%% ---------------------------------------------------------------------------
\section{Introduction}

A genetic programming system that evolves GPU programs has to deal with
illegality. A random GPU kernel can read past an array, let two threads write
one element, place a barrier under divergent control, or loop for ever. Each of
these either fails to compile, gives a nondeterministic answer, or hangs the
device. In our earlier WGSL work (Section~\ref{sec:earlier}) these failures
were handled downstream: by renderer substitutions, legality refusals and
filters. An audit of one such run counted 558 of 2,668 scorings (20.9\,\%)
refused at legality, plus ten silent substitutions in the renderer
(\src{docs/reports/2026-10-21_A1_post_expression_inventory.md}:5,13,183).

Andrew Morgan set the goal on 2026-10-21 (\src{docs/GOAL.md}): \emph{a true
multi-typed GEP kingdom for WGSL, in which every gene expresses a legal GPU
program because the symbol table's types describe exactly the legal programs.}
The chromosome, the master symbol table, the geneframe and the shared typed
GPU core are those of the core paper \cite{morgan_core}. This paper describes
what is particular to WGSL: the types that carry the rules, the generators
and witnesses that emit and prove the tables, the reader that brings existing
kernels into the kingdom, and the fitness evaluation. It then reports the first
application to a production kernel: UltraDim's k-means centroid update.

\paragraph{Contributions.}
\begin{enumerate}
\item A type discipline for WGSL in which an index in range of buffer $B$, a
  cell a thread owns, a uniform context for barriers, a bounded loop, the extent
  of a workgroup pool, a compiler-assigned scatter slot and a per-dispatch
  access mode are each a type of the master symbol table
  (Section~\ref{sec:kingdom}).
\item Generators that refuse to write a table unless witnesses over the
  emitted text prove each property, including a polynomial proof that every
  length bound is at most the interface's buffer length for every parameter of
  the dispatch contract (Section~\ref{sec:witness}).
\item A reader that brings source kernels into the kingdom exactly, so that a
  population can start from hand-written code, and a measured account of what
  the kingdom's own spelling costs in speed (Section~\ref{sec:spelling}).
\item An application to UltraDim's k-means centroid update, with the protocol,
  the controls, the results and the measurement faults found and fixed
  (Sections~\ref{sec:eval}--\ref{sec:faults}).
\end{enumerate}

%% ---------------------------------------------------------------------------
\section{Related work}

\paragraph{Typed and grammar-guided genetic programming.} Strongly typed GP
\cite{montana1995stgp} restricts each argument of a function to a type, so
that only well-typed trees are built. Grammar-guided GP \cite{whigham1995grammar}
and grammatical evolution \cite{oneill2001ge} restrict the search to the
sentences of a grammar. GEP \cite{ferreira2001gep} keeps a fixed-length linear
genotype whose head-and-tail layout guarantees that every gene expresses a
valid tree. Multi-typed GEP combines the last two: one reading frame per type
\cite{morgan_core}. What this paper adds is the use of the types for semantic
legality of GPU programs (bounds, ownership, uniformity, termination and
determinism), rather than only syntactic or type-correctness.

\paragraph{Superoptimisation and genetic improvement.} Massalin's
superoptimiser \cite{massalin1987superopt} enumerated short programs; STOKE
\cite{schkufza2013stoke} searched x86 code stochastically, with a cost of
correctness plus performance. Genetic improvement \cite{petke2018gi} evolves
existing software, and has been applied to CUDA kernels
\cite{langdon2014cuda,langdon2017gpu}. These works check candidate correctness
by testing and, for STOKE, by a verifier. In our kingdom, legality (in-range
access, race freedom, termination) holds by construction. Functional
correctness is still judged by fitness on cases.

\paragraph{Equality saturation.} E-graphs \cite{tate2009equality,willsey2021egg}
find equivalent programs by rewriting. The runner uses the kingdom's e-graph
to make equivalent variants of the start for the pump
(Section~\ref{sec:setup}). We do not evaluate it separately here.

\paragraph{Measurement.} Mytkowicz et al.\ \cite{mytkowicz2009wrong} showed that
innocuous changes of the experimental set-up, such as memory layout, bias
performance comparisons. We met the GPU form of that fault
(Section~\ref{sec:faults}).

%% ---------------------------------------------------------------------------
\section{The kingdom}
\label{sec:kingdom}

\subsection{Genotype and expression}
A chromosome is a fixed number of Karva genes joined by a linker; each gene has
a head of length $h$ and a tail of length $t=h(n-1)+1$, with $n$ the largest
arity \cite{ferreira2001gep}. The types are columns of the gene: every
position holds one symbol per type, so a gene carries one open reading frame
per type. Expression walks from the root. Each argument slot asks for a type,
and the next symbol of that type's frame fills it (\src{docs/GOAL.md}, items
1--3; the general schema is in \cite{morgan_core}). The WGSL module root is
\code{kernel.module}. Each entry point of a module interface is a type
\code{kernel.<e>}, and its argument closure is that entry's body.

A row of the table is a symbol with an output type, ordered argument types and
a WGSL spelling with holes \texttt{\{0\}}, \texttt{\{1\}}, \ldots. The renderer only
substitutes argument texts into spellings (\src{kingdoms/wgsl_true/README.md}:5).
Table~\ref{tab:rules} lists the rules and the types that carry them.

\begin{table}[t]
\caption{Each legality rule of a WGSL kernel and the types that carry it. Sources:
\src{kingdoms/wgsl_true/README.md}:15--22, \src{docs/DESIGN_inrange_index_types.md}
sections 1--11, \src{kingdoms/wgsl_true/loops.py}.}
\label{tab:rules}
\small
\begin{tabularx}{\columnwidth}{@{}p{2.5cm}X@{}}
\toprule
Rule & Carried by\\
\midrule
index in range of $B$ & \code{u32.in_<B>}: a value $<$ \code{arrayLength(&B)}; minted only by the clamp row, the literal 0, a guard register, a guard or break-guard row, or an offset into a proved run\\
length bound & \code{u32.len_<B>}: a value $\le$ the length of $B$ (\code{arrayLength}, a uniform-field bound proved by the witness, or \code{min} with one)\\
one writer per element & ownership types (\code{u32.threadrow_<b>}, \code{u32.wgown_<j>}, block and slot phases)\\
barrier in uniform control & statement classes: L (uniform), C (conditional, no barrier), W (a write phase between barriers), O (the epilogue)\\
every loop ends & trip counts of type \code{u32.size}: a size field, a literal 1..4096, and \code{+ * / >> min max} of them, never a subtraction, a free literal or a load; fuel strata \code{_f<k>}, $k<K$\\
pool no larger than named & \code{u32.wgext}: an extent of 4..$P$ words in steps of 4; every pool bound spelled \code{WG_PW}\\
no slot chosen by a gene & \code{store.cluster_<R>(u32.cluster, u32.point)}: no position argument; the slot is the compiler's\\
no read meets a write & per-entry body types: \code{block.scatter} reaches the cluster store and no load; \code{block.start} reaches loads and no store into the scratch\\
\bottomrule
\end{tabularx}
\end{table}

\subsection{In-range indices as a type}
\label{sec:inrange}
The first version of the kingdom made every load and store legal by spelling a
clamp in the row's text, \code{B[min(i, arrayLength(&B) - 1u)]}. It was legal
by construction, but it was not free. The scatter form of the k-means update
as expressed ran 2.4--2.6\,\% slower than its source, and the clamps were the
first cause (\src{docs/DESIGN_inrange_index_types.md}, Objective; logs
\src{fbench/sus2.log}, \src{sus3.log}). Being in range of $B$ is now a type,
\code{u32.in_<B>}. A load or store of $B$ takes it and spells no clamp. The
clamp is kept as one of the rows that produce the type, so every index stays
legal whatever the gene. Its table name is \code{safe.inrange_<B>}; the design
document calls it \code{inrange.<B>_clamp}. Example rows, verbatim from
\src{kingdoms/wgsl_true/mod/kmeans_update_segmented/functions.tsv}
(output type, then argument types, then spelling):

\begin{small}
\begin{verbatim}
load.scratch          u32.index  <- u32.in_scratch
   scratch[{0}]
safe.inrange_scratch  u32.in_scratch <- u32.index
   min({0}, arrayLength(&scratch) - 1u)
stmt.guard_lt_scratch_C2_f1  stmt.C2_f1
   <- u32.index; u32.len_scratch; block.C2_f1
   { let gv = {0}; if (gv < {1}) { gi_scratch2 = gv; {2} } }
block.brk_guard_scratch_B1_f0  block.B1_f0
   <- block.C1_f0; u32.index; u32.len_scratch; block.C1_f0
   {0} { let gv = {1}; if (gv >= {2}) { break; }
         gi_scratch1 = gv; } {3}
\end{verbatim}
\end{small}

The guard register \code{gi_scratch2} is declared \code{0u}. The only text
that assigns it is the guard row's, inside \code{if (gv < L)} with $L$ of type
\code{u32.len_scratch}. A read of the register anywhere therefore holds 0 or a
value some guard proved in range: possibly stale, never out of range. Fitness
judges the value (\src{docs/DESIGN_inrange_index_types.md} sections 2 and 4).
The break-guard row carries the same proof for the common loop shape
\texttt{if (i >= L) \{ break; \}}; on reply 2's histogram entry the clamp it
removes had cost 0.8\,\% (0.0609\,ms against 0.0616\,ms; section 7,
\src{fbench/b4_hist.log}). Offset types \code{u32.upto_<L>} give an index
\code{base + c} into a run that the dispatch contract places inside $B$
(section 6).

\subsection{The workgroup pool's extent}
The interface's fixed 1,024-word workgroup pool limited occupancy. On reply 2's
scatter entry, a 260-word pool ran at 0.0252\,ms against 0.0268\,ms as written
(\src{fbench/b4_pool.log}), about 6\,\%. Each entry now takes an extent of type
\code{u32.wgext}. Every pool row spells its bound as \code{WG_PW}, never a
number, and the module declares \code{WG_PW} as the largest extent its entries
name, plus 4 words (\src{src/wgsl/true_kingdom.rs}:567, \code{prologue_for}
:589). The pool witness found a hole in the fixed pool: the split phases'
writes were guarded only by the split point, which can lie past the pool
(design section 8).

\subsection{Compiler-assigned slots for the scatter}
\label{sec:scatter}
The counting-sort form of the k-means update stores each point's id at a slot
read from data: \code{scratch[ids + starts[a] + prefix[b][a] + rank]}. That
store is injective only if other entries computed the counts and scans
correctly. A gene whose histogram body is empty is legal, and then every block
stores its first point to one cell. No type of the scatter's own gene can see
this (design section 9). Two alternatives were rejected. An \code{atomicMax}
form would have changed UltraDim's driver interface. An \code{atomicAdd}
counter per cluster hands out slots in the order threads arrive, so the
cluster's points, and hence the mean's summing order and its last bits, would
change from run to run (section 10).

The adopted design, approved by Andrew Morgan on 2026-10-31: \emph{a gene names
the cluster a value belongs to; it never names a slot.} The gene sees one row,
\code{store.cluster_scratch(u32.cluster, u32.point)}, of type
\code{block.scatter}. The compiler emits a deterministic counting-sort scaffold
around it: a per-block histogram, the prefix sums, and a stable rank in
invocation-index order. Each cluster's points therefore land in index order and
the mean sums them in A7's order. The scaffold entries take about 75\,\% of the
module's time, so evolution can still choose among 26 correct variants through
a knob of type \code{kernel.hist}
(\src{kingdoms/wgsl_true/countsort.py}:159--180):
\begin{itemize}
\item 24 variants combine the histogram, offsets and starts entries of either
  the hand-written segmented form or the seed-1 champion with a scatter
  workgroup width of 64, 128 or 256;
\item one, \code{r2}, is the layout of the model-written kernel ``reply 2''
  (Section~\ref{sec:setup});
\item one, \code{none}, is the gather form.
\end{itemize}
A variant is one terminal, spelled with its pieces byte for byte. The renderer
fills every \code{<cs:piece>} marker of a module from that module's single
knob (\src{src/wgsl/true_kingdom.rs}:550).

\subsection{Access modes}
With the slot fixed, one hole remained. A gene's \code{load.scratch} could read,
inside the scatter entry, the region the cluster store writes in the same
dispatch, which is a race. The rule is now that every scratch region is
read-only or write-only within one entry, and the entry's body type carries
it. \code{block.scatter} reaches the cluster store and no load. The mean's
\code{block.start} reaches loads and no scratch write. \code{u32.point} and
\code{u32.cluster} each have one producer, so the scatter body has nothing to
compute from the scratch (design section 11; commit d86a401c).

\subsection{The tables}
Table~\ref{tab:tables} gives the sizes of the two k-means interfaces. The fuel
budget is $K=4$, because A7's loop over $n$ needs stratum 4 for $n$ up to
$2^{20}$. f16 is dropped (\code{has_f16=False}), because UltraDim requests no
device feature (\src{docs/kmeans_update_dropin.md}, gaps 1--2;
\src{kingdoms/wgsl_true/interface.py}:190,227).

\begin{table}[t]
\caption{The k-means interfaces' tables, counted on 2026-11-01 at about 01:20Z
(\code{wc -l} less one header line; the working tree held uncommitted edits to 2 and
4 lines of these files).}
\label{tab:tables}
\small
\begin{tabular}{@{}lrrrl@{}}
\toprule
Interface & types & rows & variants & bindings\\
\midrule
\code{kmeans_update} & 631 & 24,384 & -- & 0--3\\
\code{kmeans_update_segmented} & 646 & 25,030 & 26 & 0--4\\
\bottomrule
\end{tabular}
\end{table}

\subsection{The witnesses}
\label{sec:witness}
The generator \src{kingdoms/wgsl_true/gen_tables.py} runs every witness before
it writes a table (lines 81--102). It checks the scaffold byte for byte after
writing (line 105). Each witness reads the emitted text, not the generator's
intent:
\begin{itemize}
\item \textbf{in range} (\src{inrange.py}:268--365): every subscript of a
  storage binding is \texttt{B[\{i\}]} with argument $i$ of type \code{u32.in_<B>},
  or the clamp text. Only the four allowed forms mint \code{u32.in_<B>}. A guard
  register is assigned only inside its guard. Every length bound over the
  uniform's fields is at most the interface's length of $B$ for every parameter
  of the problem's \code{safety.contract}: each variable is written as its lower
  bound plus a non-negative part, and every coefficient of
  $\mathit{length}-\mathit{bound}$ must be non-negative. No bound may wrap at
  the contract's largest parameters (lines 234--247, 280--284).
\item \textbf{pool} (\src{pool.py}:343--387): every \code{wg_pool[X]} is
  \code{min(.., WG_PW)}, the clamping helper \code{xw}, or a write under a
  condition with \code{X < WG_PW} as a conjunct, and no pool text spells its
  extent as a number.
\item \textbf{counting sort} (\src{countsort.py}:268--315): every durable
  write is an own-index store or the exact cluster store, the scaffold names
  are assigned by no row, and the \code{kernel.hist} producers are exactly the
  variants plus an identity.
\item \textbf{access modes} (\src{countsort.py}:350--372): for every
  \code{kernel.entry_*}, the closure of its argument types does not hold both a
  reader and a writer of a scratch arena.
\end{itemize}
On the table as it stood before section 11, the access-mode witness named
the scatter entry: its gene could read the scratch through 3 rows and write
it through 60 (design section 11).

\subsection{The reader and the share pass}
\label{sec:spelling}
\code{read_module} (\src{src/wgsl/true_reader.rs}) reads a WGSL module into
a tree of the kingdom. The tree is then written back into a gene. A source
kernel keeps its own guards, as guard rows, and loses only the clamps its
guards make redundant (design section 5). A7, the planted control, the
scatter module, the gather module and both model replies read exactly. These
are tests \code{ultradim_kmeans_update_kernels_read_exactly_into_their_interfaces}
and \code{the_models_kmeans_replies_read_into_the_scatter_interface}.
The share pass (\src{src/wgsl/true_share.rs}:439) binds loop-invariant,
context-free subtrees at function scope. Before it, the kingdom's expression of
A7 ran at 0.579 of A7's speed, because A7's \code{let}s were inlined into its
loop over $n$. With it, A7 as expressed runs at 0.9995 of its own text
(\src{fbench/rt_champ_vs_a7.log}:14).

%% ---------------------------------------------------------------------------
\section{Fitness evaluation}
\label{sec:eval}

\subsection{The problem}
UltraDim's A7 \code{update_centroids} (qdrant
\src{lib/projection-voting-kmeans/src/gpu_kmeans.rs}:654--715) runs one
invocation per (centroid, dimension) cell. Each invocation scans all $n$
assignments in index order and keeps a Kahan-compensated \cite{kahan1965} sum
of the matching points. It writes the mean when the count is positive and
leaves an empty cluster's centroid unchanged. Its cost is $O(kdn)$. The scatter
form (\src{problems/kmeans_update_segmented/segmented.wgsl}) is UltraDim's A6
stable counting sort, rewritten with A7's semantics, as five entries
\code{hist}, \code{offsets}, \code{starts}, \code{scatter} and \code{mean}. A
fifth binding holds a scratch arena of $n+2k+2\lceil n/256\rceil k$ words.
Its cost is $O(nd)$ (\src{docs/kmeans_update_dropin.md}).

\subsection{Cases and tolerance}
There are 8 correctness cases, 4 training and 4 held out
(\src{problems/kmeans_update_segmented/cases/constraint.toml}). They cover
$n$ from 1 to 4,000, $d$ up to 2,048, $k$ up to 300, empty clusters at the end
and at random, and a partial workgroup. The expected output is A7's own output
on the same device. The tolerance on a cell is $(n+2)\,\epsilon_{f32}$ of the
expected output's RMS, Higham's $\gamma_n$ bound for any summing order
\cite{higham2002}. Timing uses its own inputs and never a correctness case. The
timing case is $n=8{,}192$, $d=256$, $k=1{,}024$. The quality group has six
inputs per case, each its own input to HFF: the share of cells outside
tolerance; the shares off by over 1, 10 and 50\,\%; and the mean and largest
error over the RMS (\src{phylu/examples/wgsl_true_case_fit.rs}:22--27, 194).

\subsection{Staged HFF}
Fitness is Andrew Morgan's Hyperspherical Fitness Function (HFF)
\cite{morgan2026hff} with the TrueNorth angle, staged
(\src{src/runner.rs}:97). Stage 1 is TrueNorth over the quality inputs. In
stage 2, an exact individual scores the two-level angle of the quality and
speed groups, in $[0,\pi/2)$. An inexact one scores $\pi/2+\theta_q$, so no
inexact kernel, however fast, ranks above an exact one. The speed group holds
11 timing slots. Each slot is one interleaved pair against the yardstick,
entered as $r/(1+r)$ with $r$ the candidate's GPU time over the yardstick's.
Pairs are never averaged before HFF (\src{src/runner.rs}:123).

\subsection{The timing protocol}
Each rule below answers a fault of Section~\ref{sec:faults}:
\begin{itemize}
\item every timed run is one unsplit dispatch per entry, as UltraDim issues it,
  for a kernel whose bound is at most twice the yardstick's; above that, an
  exact kernel is untimed and gets no speed credit;
\item a pair interleaves the two sides run by run, alternating which goes
  first (\src{src/runner.rs}:220). Each side takes its fastest run, over at
  least 1\,s of runs, and one warm-up pair is discarded;
\item the case buffers are made once and rewritten before every run (fuller
  40ae60fc);
\item every timed kernel's output on the timing case is compared with the
  yardstick's, within a cell tolerance of $3.71\times10^{-4}$; a kernel
  outside it gets no speed credit;
\item before any run, two A/A sets of 11 pairs each must lie within
  $\pm5\,\%$;
\item a speed claim counts only when the median of 11 fresh pairs is below the
  least A/A ratio \emph{and} at least 9 of the 11 pairs are faster, a sign
  test with one-sided $p=0.033$ (\src{src/runner.rs}:204).
\end{itemize}

%% ---------------------------------------------------------------------------
\section{GPU implementation}

Table~\ref{tab:where} states where each stage runs. The population is drawn,
selected and varied on the device by the shared typed core \cite{morgan_core}.
Expression, rendering and HFF run on the host. The kernels themselves run and
are timed on the device. In this runner, unlike symbolic regression, HFF runs
on the host, because the scores are a few dozen inputs per individual and
compiling each new pipeline dominates the cost: 500--730\,ms per new kernel,
measured in the earlier ST3/ST5 runs (\src{docs/PLAN_true_gep_wgsl.md}:35).

\begin{table*}[t]
\caption{Where each stage runs. fuller paths are under \src{src/} or
\src{kingdoms/}, phylu paths under \src{phylu/}; line numbers are those of the
commits read (fuller d450a10f, phylu b79f3c38).}
\label{tab:where}
\small
\begin{tabularx}{\textwidth}{@{}lp{2.2cm}X@{}}
\toprule
Stage & Runs on & Function (file:line)\\
\midrule
Table generation and witnesses & host, offline (Python) & \src{kingdoms/wgsl_true/gen_tables.py}:81--105\\
Table load & host & \code{TrueKingdom::load}, \src{src/wgsl/true_kingdom.rs}:202\\
Initial population & device & \code{TypedDevice::init}, \src{phylu/src/evolve/typed_device.rs}:507 (called \src{true_wgsl_fit.rs}:771)\\
Selection, mutation, inversion, transposition, crossover & device & \code{profile_vary} \src{typed_device.rs}:665; \code{vary_parents} :747 (called \src{true_wgsl_fit.rs}:1445, 2106)\\
Behaviour-diverse parents (\code{--diverse}) & host & \code{diverse_parents}, \src{phylu/src/evolve/true_wgsl_fit.rs}:1965\\
Model edit (Sonnet, \code{--ops model=P}) & host, external process & \code{call_model}, \src{phylu/examples/wgsl_true_case_fit.rs}:1208\\
Karva expression of a gene & host & \code{TrueModule::tree}, \src{src/wgsl/true_module.rs}:92; \src{phylu/src/evolve/true_wgsl_fit.rs}:1741\\
Share pass and rendering to WGSL & host & \code{Share::share} \src{src/wgsl/true_share.rs}:439; \code{TrueKingdom::render} \src{true_kingdom.rs}:516, \code{prologue_for} :589\\
Reading a source kernel & host & \code{read_module}, \src{src/wgsl/true_reader.rs}\\
E-graph variants of the start (for the pump) & host & \src{phylu/examples/wgsl_true_case_fit.rs}:964\\
naga parse and validation & host & \src{src/wgsl/naga_shape.rs}:37; \src{src/wgsl/target.rs}:93\\
Pipeline compile & host driver, for the device & \code{DeviceExecutor::compile}, \src{src/wgsl/harness.rs}:3517\\
Work bound and dispatch split & host decides; device runs & \code{tree_work_bound_workgroup_in} \src{true_kingdom.rs}:329; \src{work_guard.rs}:130; \src{harness.rs}:3313--3350\\
Correctness cases & device & \code{CaseSet::judge} \src{src/wgsl/case_set.rs}:130; \code{DeviceExecutor::run} \src{harness.rs}:3605\\
Output comparison, quality inputs & host & \src{case_set.rs}:139--140; \src{phylu/examples/wgsl_true_case_fit.rs}:194\\
Timing & device timestamps; host schedules & \src{harness.rs}:3361, 3386; \code{time_unsplit} \src{case_set.rs}:193; \src{src/runner.rs}:220; \src{wgsl_true_case_fit.rs}:582--594\\
Timing-case output check & host & \src{phylu/examples/wgsl_true_case_fit.rs}:617\\
Staged HFF, verdict & host & \code{staged_hff} \src{src/runner.rs}:97; \code{speed_verdict} :204\\
\bottomrule
\end{tabularx}
\end{table*}

\paragraph{The work guard.} After a run of random genes at $K=2$ starved
the macOS window server and forced a hard reset on 2026-10-21 (fuller
9f770a28), every dispatch runs under a work guard. The guard bounds each
submit's node executions from the tree's own work bound, and splits a dispatch
into submits of whole workgroups (\src{src/wgsl/work_guard.rs}:24, 130). A7's
own timing gives a measured rate of $1.22\times10^{11}$ node executions a
second. From it the measured cap is $1.22\times10^{9}$ per submit, a tenth of a
second over a safety factor of 10 (:49--54). The splits serve the correctness
cases only. Speed is timed unsplit.

\paragraph{Determinism.} A CPU executor that can run workgroups in reverse
order (\src{src/wgsl/cpu.rs}:878) tests the kingdom's determinism claim.
300 random modules per k-means table give the same bits with their workgroups
reversed (test
\code{random_kmeans_modules_give_the_same_bits_with_their_workgroups_reversed}).
The CPU executor is a test reference only.

%% ---------------------------------------------------------------------------
\section{Experiments}
\label{sec:exp}

\subsection{Setup}
\label{sec:setup}
All runs used an Apple M3 Max under Metal through wgpu 30.0.1, the runner
\src{phylu/examples/wgsl_true_case_fit.rs}, regime P (from given kernels toward
faster), a population of 24 and a budget of 1,400 fitness evaluations. The
start was A7 in the five-entry interface (\src{gather.wgsl}). The hand-written
scatter form was added to generation 0 with \code{--also}. The rest of
generation 0 was 22 point mutants of the start, each with a text and a tree
different from every member's. The start's e-graph variants were fed to the
pump. Operators were point mutation, inversion, transposition, crossover and
marriage (\src{fbench/kmeans_P_s1/population.tsv},
\src{per_generation.tsv}). ``Reply 2'' is a kernel that the model-edit operator
(Sonnet) wrote in a smoke run. On that day it did not read into the kingdom
(\src{fbench/kmeans_model_smoke2/model_calls.tsv}). Its text is now a reader
test and the source of the \code{r2} variant
(\src{tests/corpus/kmeans_model_replies/}).

\paragraph{Expectations, stated before the runs.} The scatter form's
$O(nd)$ cost against A7's $O(kdn)$ should give most of any gain. Evolution was
expected to find small gains in the scaffold's widths and the mean's loop. The
kingdom's expression of a kernel was expected to run at parity with the kernel
as written once the in-range types were built (design section 10, target
parity).

\subsection{Results}
Table~\ref{tab:retime} gives the re-time of the champion of run
\code{kmeans_P_s1} after the buffer fix (2026-10-31, 14:18--14:42Z). Every
kernel in the table gave A7's output bit for bit on all 8 cases (0 cells
differing), and on the timing case 0.00\,\% of cells lay outside the
yardstick's tolerance. All six A/A sets passed.

\begin{table}[t]
\caption{Interleaved re-times, 11 pairs each, timing case $n$ 8,192, $d$ 256,
$k$ 1,024. Ratio is candidate over yardstick (lower is faster). Logs:
\src{fbench/rt_champ_vs_a7.log}, \src{rt_champ_vs_scatter.log},
\src{rt2_scatter_vs_a7.log}, \src{rt2_reply2_vs_scatter.log}.}
\label{tab:retime}
\small
\begin{tabular}{@{}lrrr@{}}
\toprule
Candidate / yardstick & median & faster & A/A range\\
\midrule
champion / A7 & 0.0274 & 11/11 & 0.9998--1.0002\\
champion / scatter form & 0.9839 & 11/11 & 0.9961--1.0055\\
scatter form / A7 & 0.0278 & 11/11 & 0.9662--1.0086\\
reply 2 / scatter form & 0.9925 & 11/11 & 0.9970--1.0020\\
\midrule
\multicolumn{4}{@{}l}{\emph{the kingdom's expression, against the same text as written}}\\
A7 & 0.9995 & -- & \\
scatter form & 1.0255 & 0/11 & \\
scatter form (second run) & 1.0238 & 0/11 & \\
\bottomrule
\end{tabular}
\end{table}

The absolute times on the timing case were 8.452\,ms for A7 and about
0.234\,ms for the scatter form, a whole module of five dispatches
(\src{rt_champ_vs_a7.log}:14, \src{rt_champ_vs_scatter.log}:15). Per entry,
the champion's own text took 0.0509, 0.0119, 0.1095, 0.0278 and 0.0290\,ms for
\code{hist}, \code{offsets}, \code{starts}, \code{scatter} and \code{mean}
(\src{rt_champ_vs_a7.log}:17--18).

\paragraph{What evolution changed}
(\src{docs/kmeans_champion_s1/README.md}):
\begin{enumerate}
\item In \code{mean}, the Kahan subtraction now reads a register that is never
  written. The sum is therefore a plain sum. The compensation is still
  computed, but nothing reads it.
\item \code{hist} runs at workgroup width 128 instead of 256.
\item Three rewrites have no effect: a loop's final counter in place of the
  literal 256, knob constants that no line reads, and a dead store.
\end{enumerate}

\paragraph{The Kahan point.} On Metal, A7 with its Kahan lines replaced by a
plain sum in the same order gives A7's bits on 8 of 8 cases, with 0 cells
differing. Reversing the order or summing in blocks of 64 changes 27,832 and
23,026 cells, so the cases are large enough to show a working compensation
(\src{docs/reports/2026-10-28_metal_removes_a7_kahan.md}). The likely cause is
the Metal compiler folding $(t-s)-y$ under reassociation. We did not inspect
the compiled shader, so this is an inference. The champion's main change is
therefore exact on Metal. It need not be exact on Vulkan, which has not been
measured.

\paragraph{Learning curve.} Figure~\ref{fig:curve} shows the best exact
kernel's speed ratio by generation, as measured inside the run (one slot,
before the buffer fix). Generation 0's best was the scatter form at 0.0285 of
A7. The run reached 0.0276 at generation 9 and 0.0273 at generation 13, and
did not improve in the remaining 84 generations. It stopped at 1,415
evaluations after 36.4 minutes, with 5 work-guard failures: mutants with a
single workgroup over the per-submit cap (\src{fbench/kmeans_P_s1/report.tsv},
\src{bugs.log}).

\begin{figure}[t]
\centering
\Description{Two step plots against generation: the best exact kernel's speed ratio against A7 falls from 0.0285 to 0.0276 at generation 9 and 0.0273 at generation 13, then stays flat; the best HFF falls in the same steps.}
\pgfplotstableread[col sep=tab]{wgsl_multityped_gep_data/kmeans_P_s1_per_generation.tsv}\curve
\begin{tikzpicture}
\begin{axis}[width=\columnwidth, height=4.4cm, xlabel={generation},
  ylabel={best exact / A7}, ymin=0.0270, ymax=0.0288,
  yticklabel style={/pgf/number format/fixed, /pgf/number format/precision=4},
  scaled y ticks=false, xmin=0, xmax=96, tick label style={font=\footnotesize},
  label style={font=\footnotesize}]
\addplot[const plot, thick, blue!70!black] table[x=generation, y=best_exact_speed_ratio]{\curve};
\end{axis}
\end{tikzpicture}
\begin{tikzpicture}
\begin{axis}[width=\columnwidth, height=4.0cm, xlabel={generation},
  ylabel={best HFF}, xmin=0, xmax=96,
  yticklabel style={/pgf/number format/fixed, /pgf/number format/precision=4},
  scaled y ticks=false, tick label style={font=\footnotesize},
  label style={font=\footnotesize}]
\addplot[const plot, thick, red!70!black] table[x=generation, y=best_hff]{\curve};
\end{axis}
\end{tikzpicture}
\caption{Run \code{kmeans_P_s1}: the best exact kernel's speed ratio against A7
(in-run, one slot) and the best HFF (0 is best), by generation. Data:
\src{wgsl_multityped_gep_data/kmeans_P_s1_per_generation.tsv}, from the run's
\src{per_generation.tsv}.}
\label{fig:curve}
\end{figure}

\begin{figure}[t]
\centering
\Description{Scatter of 11 pair ratios: both A/A sets lie between 0.996 and 1.006, the champion's pairs between 0.981 and 0.987, and the scatter form as then expressed between 1.022 and 1.028.}
\pgfplotstableread[col sep=tab]{wgsl_multityped_gep_data/champion_vs_scatter_pairs.tsv}\pairs
\begin{tikzpicture}
\begin{axis}[width=\columnwidth, height=4.6cm, xlabel={pair},
  ylabel={ratio to scatter form}, xmin=0.5, xmax=11.5, ymin=0.975, ymax=1.032,
  xtick={1,...,11}, legend columns=2, legend style={font=\scriptsize, at={(0.5,-0.32)}, anchor=north},
  yticklabel style={/pgf/number format/fixed, /pgf/number format/precision=3},
  tick label style={font=\footnotesize}, label style={font=\footnotesize}]
\addplot[only marks, mark=o, gray] table[x=pair, y=aa1]{\pairs};
\addplot[only marks, mark=x, gray] table[x=pair, y=aa2]{\pairs};
\addplot[only marks, mark=*, blue!70!black] table[x=pair, y=champion_text]{\pairs};
\addplot[only marks, mark=square*, red!70!black] table[x=pair, y=scatter_expressed]{\pairs};
\legend{A/A set 1, A/A set 2, champion, scatter form as expressed}
\end{axis}
\end{tikzpicture}
\caption{The 11 pairs against the hand-written scatter form, with both A/A sets
and the scatter form as the kingdom then expressed it. Data:
\src{wgsl_multityped_gep_data/champion_vs_scatter_pairs.tsv}, from
\src{fbench/rt_champ_vs_scatter.log}.}
\label{fig:pairs}
\end{figure}

\paragraph{The cost of the kingdom's spelling, after sections 1--11.}
Table~\ref{tab:perentry} is the per-entry pass that tested the target of
parity, made after the in-range types, the pool extent, the compiler-assigned
slots and the access modes were built. The segmented form as expressed is at
parity: 0.0001\,ms over its text in total, against an A/A floor of 0.0003\,ms.
Reply 2 as expressed is 0.0003\,ms over its text in total, within the floor.
Its scatter entry, however, is 0.0003\,ms (1.1\,\%) slower in both passes,
where that entry's A/A floor is 0.0000. The cause is the scaffold's \code{r2}
prologue. It computes the slot for every invocation and stores under a flag,
where reply 2 returns early, and the cluster store clamps its slot (design
section 10).

\begin{table}[t]
\caption{Per entry, ms, the fastest of 4,302 round-robin runs over 4 pipelines
a text, 2026-10-31 23:33Z (\src{fbench/b7b_pe.log}). ``Yardstick'' is the A/A
reference, the segmented form's text.}
\label{tab:perentry}
\small
\setlength{\tabcolsep}{3pt}
\begin{tabular}{@{}lrrrrrr@{}}
\toprule
Text & hist & offs. & starts & scat. & mean & total\\
\midrule
yardstick & 0.0555 & 0.0118 & 0.1064 & 0.0290 & 0.0287 & 0.2315\\
segmented, written & 0.0560 & 0.0118 & 0.1064 & 0.0290 & 0.0287 & 0.2318\\
segmented, expressed & 0.0560 & 0.0118 & 0.1064 & 0.0291 & 0.0287 & 0.2319\\
reply 2, written & 0.0613 & 0.0117 & 0.1020 & 0.0265 & 0.0285 & 0.2301\\
reply 2, expressed & 0.0614 & 0.0117 & 0.1020 & 0.0268 & 0.0285 & 0.2304\\
\bottomrule
\end{tabular}
\end{table}

\paragraph{The other runs.} Table~\ref{tab:runs} lists every k-means P run of
2026-10-31. Two champions were withdrawn because they did no work on the timing
case. The timing-case output check was added because of them.

\begin{table}[t]
\caption{The k-means P runs of 2026-10-31 (\src{fbench/kmeans_P_progress.tsv};
checks \src{fbench/check_kmeans_*.log}, 11:06--11:14Z, before the buffer fix).
Ratios are against A7.}
\label{tab:runs}
\small
\setlength{\tabcolsep}{3pt}
\begin{tabularx}{\columnwidth}{@{}lrrX@{}}
\toprule
Run (\code{kmeans_P_}\ldots) & evals & in run & status\\
\midrule
\code{s1} & 1,415 & 0.0273 & valid; re-time 0.0274\\
\code{s2} & 348 & 0.0252 & stopped; never checked on the timing case\\
\code{diverse_s2} & 1,400 & 0.0284 & check 0.0284, 0\,\% outside\\
\code{model_s3} & 1,129 & 0.0285 & check 0.0282, 0\,\% outside\\
\code{long_s2} & 4,013 & 0.0260 & withdrawn: 99.88\,\% of timing-case cells wrong\\
\code{planted_diverse_s1}$^\dagger$ & 1,402 & 0.0022 & withdrawn: no work on the timing case\\
\bottomrule
\multicolumn{4}{@{}l}{\footnotesize $^\dagger$ on the one-entry interface \code{kmeans_update}, from the planted kernel.}
\end{tabularx}
\end{table}

%% ---------------------------------------------------------------------------
\section{Faults found and fixed}
\label{sec:faults}

Each fault in Table~\ref{tab:faults} produced a result that looked better
than it was, or hid a cost. Each was found before the numbers of
Table~\ref{tab:retime}, and each fix is now part of the protocol.

\begin{table*}[t]
\caption{Faults in our own method, found and fixed. Sources:
\src{docs/papers/2026-10-31_kmeans_centroid_update.md} section 3,
\src{docs/kmeans_update_dropin.md} (timing rules), and the logs named.}
\label{tab:faults}
\small
\begin{tabularx}{\textwidth}{@{}p{3.2cm}XX@{}}
\toprule
Fault & What it did & Fix\\
\midrule
Split timing & Timing a dispatch split into many submits favoured kernels with small work bounds, and was not UltraDim's dispatch & every timed run is one unsplit dispatch per entry\\
Display noise & The display shares the GPU; self-timing ranged 0.92--1.10 when each side was timed in its own block & run-by-run interleaving, fastest of at least 1\,s a side, one warm-up pair discarded\\
Buffer placement & Fresh buffers every run: speed followed where they landed; one entry moved by up to 25\,\%, one text compiled separately ran at 0.64--0.90\,ms & case buffers made once, rewritten before every run (fuller 40ae60fc); five compilations of one text then agree within 0.2\,\%\\
The u32-overflow cheat & A guard \code{gid.x >= (k*d)^2} wrapped to 0 at the timing size, so the kernel did no work there and claimed up to 450$\times$ & every timed kernel's timing-case output is checked against the yardstick's; two champions withdrawn (Table~\ref{tab:runs})\\
The expression cost & The kingdom's spelling cost speed: A7 as expressed ran at 0.579 of its text (inlined \code{let}s in the loop); the scatter form at 1.024--1.026 (clamps, the pool, the scatter's slot) & the share pass; in-range, offset and break-guard types; the pool extent; compiler-assigned slots. A7 now 0.9995, the segmented form at parity (Table~\ref{tab:perentry})\\
The dispatch limit & A mutated workgroup width needed 262,144 workgroups (the limit is 65,535) & the types bound every workgroup width; the empty entry is still open (Section~\ref{sec:limits})\\
Unbounded loops & 6 smoke-run mutants had a work bound of 4,218,022,233 or a bound read from \code{assignments}, which the guard could not split & a trip count is of type \code{u32.size}: sizes, literals, and \code{+ * / >> min max} of them\\
\bottomrule
\end{tabularx}
\end{table*}

%% ---------------------------------------------------------------------------
\section{Earlier WGSL work}
\label{sec:earlier}

The kingdom grew out of four earlier lines of work. Each one taught the
method something.
\begin{itemize}
\item \textbf{Weighted Gram superoptimisation (2026-10-10 to 14).} A rule
  beam over 24 correct, model-written kernels gave 1.25$\times$ over generation
  0's best (6 of 6 paired rounds; 247 of 247 rule children correct). Random
  starts never produced a correct kernel. Single-round timing had spread
  champions from 6.0$\times$ to 8.7$\times$ when re-timing put them all at about
  6.1$\times$. 37 of 40 random kernels raced, and naga's raw \texttt{\%} hung the GPU
  on modulo by zero (\src{docs/reports/2026-10-13_SUMMARY.md};
  \src{docs/PLAN_hgp_gram.md}:193--217).
\item \textbf{V24 streaming index (2026-10-17 to 19).} The best individual was
  a quantised exhaustive scan at recall@10 of 1.0 up to 15,000 records. At that
  size the fitness could not tell a scan from an index
  (\src{docs/reports/2026-10-17_v24_overnight.md}:3--5).
\item \textbf{cuneus image shaders (ST9--ST13).} volumepassage and nebula were
  read exactly (0 of 1,048,576 channels differ), and the expressed volumepassage
  reached 0.983 of the shipped shader's time. A planted slowdown was recovered
  (1.742$\times$ to 0.962$\times$), but no faster shader was found. An apparent
  2.2\,\% gain existed only under tiled submits, and the fitness allowed a literal
  to be replaced by a uniform that never varied across cases
  (\src{docs/reports/2026-10-22_ST9_volumepassage_champion/README.md};
  \src{docs/reports/2026-10-23_summary_instance_and_class.md}:125--129).
\item \textbf{UltraDim ingest (UI1--UI4).} 50 of 67 ingest kernels read
  exactly. The best evolved kernel was 59--107$\times$ faster at ingest but had
  lower recall (0.62 to 0.38, against 0.98 to 0.97), and one champion
  hard-coded the fitness case's query count
  (\src{docs/reports/2026-10-23_UI_ingest_overnight.md}:53--71).
\end{itemize}
What this taught: legality belongs in the types (A1--A3), timing needs a noise
floor and a sign test, and a fitness must check what the timed kernel actually
computed. The k-means protocol is built on these lessons.

%% ---------------------------------------------------------------------------
\section{Discussion}

Legality by construction worked as designed on this problem. Every gene
expresses a kernel that compiles and runs. 300 random modules per table
validate under naga, index in range on the CPU, never exceed the dispatch
limit, and give the same bits with their workgroups reversed (tests in
\src{src/wgsl/true_reader.rs}). The failures that remained in the run were
the work guard's: 5 of 1,415 evaluations.

The speed result needs to be read in three parts. First, the factor of about 36
over A7 is the algorithm. The hand-written counting sort is at 0.0278, and
evolution did not invent it: the scatter form was given to generation 0.
Second, evolution's own gain is 1.6\,\% over that form, and it held under the
strict protocol (11 of 11 pairs, outside both A/A ranges). Third, most of that
gain is the removal of a compensation that the Metal compiler appears to remove
already. The change is free on Metal for that reason, and it may not be free
elsewhere.

The kingdom's spelling had been costing as much as evolution gained: 2.4--2.6\,\%
on the scatter form. That cost is now removed on the segmented form, by types
rather than by special cases in the renderer. This matters for regime P,
because otherwise evolution must first win back the kingdom's own overhead
before it can improve on the source.

%% ---------------------------------------------------------------------------
\section{Limitations and open work}
\label{sec:limits}

\begin{itemize}
\item \textbf{The empty mean exceeds the dispatch limit.} The empty entry
  (\code{own.kernel_<e>}) was spelled at workgroup width 1. For the mean this
  gives $\lceil kd/1\rceil=262{,}144$ workgroups on the timing case, over
  65,535. At the time of writing, a fix is in fuller's working tree,
  uncommitted: an \code{empty_width} rule and a \code{dispatch_witness} in
  \src{kingdoms/wgsl_true/skeleton.py}.
\item \textbf{Reply 2's scaffold is 1.1\,\% slow in its scatter entry}
  (Table~\ref{tab:perentry}). An uncommitted change of the \code{r2} prologue
  returns early as reply 2 does and drops the slot clamp. It has not been
  timed. While that edit was in progress, the committed segmented table
  failed the counting-sort witness against the edited generator. It must be
  regenerated.
\item \textbf{Data-dependent indices} (a slot read from the scratch, a label
  read from \code{assignments}) still read through the clamp row (design
  section 12).
\item \textbf{One shape, one machine, one backend}: $n$ 8,192, $d$ 256, $k$
  1,024, on an M3 Max under Metal through wgpu 30. Larger $n$, rolling
  k-means' small batches, and Vulkan are not measured. The Kahan point is
  Metal-only.
\item \textbf{One valid seed.} Only \code{kmeans_P_s1} was re-timed after the
  buffer fix. Two other runs passed the timing-case check before it, and one
  stopped run was never checked (Table~\ref{tab:runs}).
\item \textbf{The timing-case check is a tolerance}, not a bit count: 0.00\,\%
  of cells outside $3.71\times10^{-4}$ of the yardstick's. Bit-exactness is
  measured on the 8 correctness cases.
\item \textbf{Not yet in UltraDim.} The change request
  (\src{docs/papers/2026-10-31_kmeans_centroid_update.md}) asks the UltraDim
  team to validate the kernel in UltraDim, end to end, on Metal and Vulkan.
  We changed no UltraDim code.
\end{itemize}

%% ---------------------------------------------------------------------------
\section{Conclusion}

A WGSL kingdom in which every gene expresses a legal kernel can be built by
putting each legality rule into a type of the master symbol table, and having
generators prove the emitted tables. Applied to UltraDim's k-means centroid
update with UltraDim's own bindings, the kingdom read the production and
hand-written kernels exactly, expressed them at parity with their text, and
evolved a kernel at 0.0274 of A7's time, bit for bit on every correctness case.
Evolution's own gain over the hand-written form is 1.6\,\%, and its main source
is specific to Metal. The open work is the two bugs above, measurement on
Vulkan and other shapes, and validation inside UltraDim.

\begin{acks}
The design of the multi-typed GEP schema, the geneframe, HFF and TrueNorth, the
pump, and the rules of the kingdom and of the timing protocol are Andrew
Morgan's.
\end{acks}

\bibliographystyle{ACM-Reference-Format}
\bibliography{wgsl_multityped_gep}

\end{document}
